% !TeX root = ../papers/nc-p3.tex
\section{固定曲線の反対階数と階数14の楕円曲線}
\label{sec:fixed-opposite-rank-drop}
\begin{theorem}\label{thm:fixed-opposite-rank-drop}
標数零の全通常倍率で、負の固定曲線の許容窓の反対切断は可逆である。
正の固定曲線では $A,B$ の第一行をそれぞれ $(a,b),(e,f)$ とすると、反対階数は
\[
 16-\boldsymbol{1}_{a^2+e^2=2zbf}-\boldsymbol{1}_{z(b^2+f^2)=2ae}.
\]
階数14,15の全窓は全始点で長さ4,5の標準次元を持つ。
階数14の窓が存在する必要十分条件は
\[
             (\lambda+1)^2(\lambda z+1)^2=16\lambda^2 z .
\]
この曲線上の各通常倍率で階数14の幾何点はちょうど四つある。
標数零の代数閉体上で、倍率も動かしたこの軌跡は二つの成分を持ち、
それぞれの滑らかな完備化は $Y^2=2(t-1)(t+1)(t+3)$ と同型で、$j=1728$ である。
\end{theorem}
\begin{proof}
反対切断を四文字のブロックに分ける。$b=1$、$\Delta=a^2-z$、
$(e,f)=u(N,D)$ と書く。文字2,3の行列式は両符号とも
$-(\lambda-1)^2(\lambda z-1)^2/\lambda^2$ である。
負の二つの残りの分子は、通常倍率の単元を除き
$(a^2-z^2u^2)^2,(a^2u^2-1)^2$ である。
曲線の二次式 $G$ と各因子の終結式を取ると、通常倍率の単元と
$a,\Delta$ の因子だけが残る。負の座標端は非許容である。

正の文字0では最後の行を $z$ で割り、文字1では二列目を $z$ で割る。
二つのブロックはともに
\[
 \begin{pmatrix}m&n&1&1\\ n&m&1&1\\1&1&C&D'\\1&1&D'&C\end{pmatrix}
\]
となる。文字0で $(m,n)=(a,e)$、文字1で $(m,n)=(b,f)$ である。
反対称部分の二つの係数は $m-n,\ k(m-n)$ で、$b=1$ では
$k=-1/(\lambda u z\Delta),-1/(\lambda u\Delta)$ である。
$m-n$ と $G$ の終結式の因子は定理\ref{thm:uniform-fixed-quartic}の
座標端とその交換像、$\Delta$、通常倍率の単元に限られる。
端点の公式では $a-e,b-f$ の分子は通常倍率の単元であり、行列交換でも非消滅を保つ。
従って反対称部分は可逆である。対称部分の二次の行列には非零の成分2があり、
その行列式は順に $-k(a^2+e^2-2zbf)$ と $-k(b^2+f^2-2ae/z)$ になる。
斉次表示で $\Delta=a^2-zb^2$ に置き換えれば、これらの恒等式は座標端にも延びる。
古い四次障害の正の四表示に二つの条件式を代入すると、いずれも通常倍率の単元である。
反対行列式の四つの障害点も定理\ref{thm:opposite-fixed-dets}で反対階数16である。
従って定理\ref{thm:uniform-fixed-quintics}は階数14,15の全点に標準次元を与える。
以上のブロックと終結式は \path{computations/verify_varying_opposite_rank.py} が検算する。

両条件が成り立つのは $(a+e)^2=z(b+f)^2$ かつ
$(a-e)^2=-z(b-f)^2$ のときである。$a+e=b+f=0$ は
$B=-A$ となり、固定点のない二位置移動に反する。従って $r^2=z$ の符号を
$a+e=r(b+f)$ により選べる。$i^2=-1$ のどちらかを取って
$a-e=ir(b-f)$ とし、$X=a/r$, $w=\lambda r$ と置く。この二式から
\[
 f=F=(1+i)X-ib,\qquad e=rE=r\big((1-i)b+iX\big).
\]
交換の比 $[e:f]=[N:D]$ は次の二次形式を与える。
\[
 K=E\big((w^2+1)X-2wb\big)-F\big(2wX-(w^2+1)b\big).
\]
$T=rw^2+wr^2-4rw+w+r$ とする。二つの補完条件
$f=f_0+f_1(a^2-zf^2)$、$e=e_0+e_1(e^2-zb^2)$ は、この代入後に共通分母
$w(r^2+1)(X+b)$ を持ち、分子は二次形式となる。分母が零になる許容点はない。
\path{computations/verify_fixed_rank_fourteen.py} は、両分子と $K$ の二変数斉次終結式が
$(w+1)^2T^2$ と、それぞれ
\[
  -2i\big(w^2+(-2+2i)w+1\big),\qquad
    2ir^2\big(w^2+(-2-2i)w+1\big)
\]
との積であることを検算する。括弧内の差は $4iw\ne0$ なので $T=0$ が必要である。
逆に $T=0$ では各分子と $K$ の係数行の全二次小行列式が消える。
$K$ の判別式は $16w(w-1)^2\ne0$ であり、二つの根は全補完条件を満たす。
実際、正の固定直線では $g_1=zf_1,h_1=e_1$ で残りの線型式も従い、
交換の比が $zf_1f+e_1e=0$、従って二つの二次式の導関数値の和が零となる対根を選ぶ。

許容性も全ての倍率で確認できる。$K$ と $X^2-b^2,E^2-F^2$ の終結式はともに
$-4(w^2-1)^2$、$X^2-F^2,E^2-b^2$ との終結式はともに
$-4(w-1)^2(w^2-6w+1)$ である。$T=r(w^2-6w+1)+w(r+1)^2$ なので、
最後の因子が零なら $r=-1$ となり通常性に反する。これらは斉次終結式なので無限遠の根も含む。
$T=0$ は $(\lambda+1)(\lambda z+1)=4\lambda r$ と同値であり、
倍率から $r$ が一意に定まる。二つの $i$ に共通する点は
$a=e,b=f$ となるため存在しない。これで各倍率の四点を得る。

倍率の曲線は
\[
 r=\frac{-2(t+3)}{t^2-1},\qquad \lambda=-\frac{(t-1)^2}{(t+3)^2}
\]
で有理に表示できる。逆は $t=(2r^2\lambda+r^2-5r+2)/(r(r-1))$ である。
各 $i$ で二次式 $K$ の根を添加することは、その判別式から $\sqrt w$ を添加することと等しい。
$w=2(t-1)/((t+1)(t+3))$ は平方でなく、$Y=(t+1)(t+3)\sqrt w$ が主張の滑らかな三次式を満たす。
$t+1$ を変数に取ると右辺は $2((t+1)^3-4(t+1))$ であり、$j=1728$ が従う。
\end{proof}

\begin{proposition}\label{prop:fixed-rank-drop-regularity}
倍率 $(3,2)$ で正の曲線の反対階数が16を下回るのは八幾何点で、全て階数15の標準Koszul左右AS正則窓である。
$(36,-1/4)$ では正の曲線の階数低下点は四幾何点で、全て階数14の標準Koszul左右AS正則窓である。
前者は倍率平面の非空開近傍で、後者は上の倍率曲線上の非空開近傍で成立する。
従って階数14の各楕円成分で全次数の正則性が未確定な許容点は高々有限個である。
\end{proposition}
\begin{proof}
正の二つの行列式を $G$ と消去し、許容な座標端を別に検算する。
$(3,2)$ の残余第一座標の多項式は文字0,1の順に
\[
 932a^4-8229a^3+18228a^2-5265a-576,\quad
 64a^4+1755a^3-18228a^2+24687a-8388 .
\]
$(36,-1/4)$ では両者とも $769a^4+3744a^3-10296a^2+134784a+996624$ となる。
各四次式は既約で、その根の上の該当点は $u=U(a)$ と一意に定まる。
後者では二つのブロックが同じ $U$ を与える。
\path{computations/verify_fixed_opposite_rank_drop.py} はこれらの等式と二座標端の可逆性を確かめ、
各閉点で101を上に持つ単純な剰余根を一つ取る。
その $(a,u)$ は順に $(40,4),(33,8),(15,71)$、反対階数は $15,15,14$ である。
消えるブロックが与える上界と剰余階数が一致する。
全ての剰余点で四次中心の文字次元は $(3,2,0,0)$、四式で割った各始点の商は
命題\ref{prop:fixed-residual-regularity}と同じ次数13までの次元を持つ。
同じ中心の持上げと自由基底の証明で、全共役の全次数の正則性が従う。

倍率を動かした終結式の残余は両文字で各四次式一つで、いずれも重複度1である。
$(3,2)$ への特殊化は次数を保ち、二式は分離的かつ互いに素である。
従って開近傍で有限エタールな八根を持つ。各根の上で共通の $u$ を与える一次の
部分終結式の最高次係数と、許容性および座標端の行列式はこの近傍を縮めれば零にならない。
$(36,-1/4)$ では倍率曲線上の四根が前定理で有限エタールであり、各四次式が全根を占める。
両場合とも有限被覆で正則性の証明に使った小行列式の非消滅が開であるから結論を得る。
\end{proof}
